\documentclass[12pt]{article}
\usepackage[utf8]{inputenc}
\usepackage{graphicx}
\usepackage{hyperref}
\usepackage{amssymb}
\usepackage{amsmath}
\usepackage[T1]{fontenc}
\usepackage{lmodern}
\usepackage{textcomp}
\usepackage{pgfplots}
\usepackage{pgfplotstable}
\usepackage{longtable}
\usepackage[capitalise]{cleveref}
\usepackage{caption}
\pgfplotsset{compat=1.18}
\captionsetup[table]{position=top,justification=raggedright,singlelinecheck=false}
\captionsetup[longtable]{position=top,justification=raggedright,singlelinecheck=false}
\captionsetup[figure]{justification=raggedright,singlelinecheck=false}
\setlength\LTleft{0pt}
\setlength\LTright{0pt}
\setlength\LTcapwidth{\textwidth}
\pgfplotstableset{
  show variable once/.style={
    columns/variable/.style={
      string type,
      column name=Variable,
      postproc cell content/.code={
        \pgfkeysgetvalue{/pgfplots/table/@cell content}{\cellcontent}
        \edef\currentvar{\cellcontent}
        \ifx\currentvar\lastvar
          \pgfkeyssetvalue{/pgfplots/table/@cell content}{}
        \else
          \ifnum\pgfplotstablerow>0
            \pgfkeyssetevalue{/pgfplots/table/@cell content}{\noexpand\hline\space\cellcontent}
          \fi
          \xdef\lastvar{\currentvar}
        \fi
      }
    }
  }
}
\def\lastvar{}
\setlength{\parskip}{0.8em}
\setlength{\parindent}{0pt}
%\renewcommand{\thesubsection}{\thesection.\alph{subsection}}

\title{\Large\bfseries ASSIGNMENT 3: GENERAL LINEAR MODEL — FITTING
AND INFERENCE}
\author{Chenghao Wang}
\date{\today}

\begin{document}

\maketitle

\section{ Model Specification}

\subsection{Model formulation}

The model is specified in \cref{subsec:covariance-structure}.

\subsection{Visualization}

The observed data along with the fitted mean structure
is demonstrated in \cref{subsubsec:fitted-mean-trajectories}.

The model included planned visit time (unit: month)
instead of actual visit time (unit: day). As shown in \cref{fig:measurement-timing-deviation},
such design appears plausible.

As is illustrated in \cref{fig:mean-trajectories-by-treatment},
the mean PostBD FEV1 trajectories of each treatment group
over visit time
appears to be linear, except at a few time points.
Actually, PostBD FEV1 exhibits high missingness 
at those time points.
Those missingness appears to introduce considerable bias.
Thus we filtered out time points with $\ge 75\%$ missing values in PostBD FEV1.
The filtered data shows a linear trend between mean PostBD FEV1
and visit time (\cref{fig:filtered-mean-trajectories-by-treatment}).



\begin{figure}[htbp]
  \centering
  \includegraphics[width=0.9\textwidth]{../hw3_code/output/measurement_timing_deviation.png}
  \caption{Measurement timing deviation across visits.}
  \label{fig:measurement-timing-deviation}
\end{figure}

\begin{figure}[htbp]
  \centering
  \includegraphics[width=0.9\textwidth]{../hw3_code/output/mean_trajectories_by_TG_over_visitc.png}
  \caption{Mean trajectories by treatment group over visits.}
  \label{fig:mean-trajectories-by-treatment}
\end{figure}

\begin{figure}[htbp]
  \centering
  \includegraphics[width=0.9\textwidth]{../hw3_code/output/filtered_mean_trajectories_by_TG_over_visitc.png}
  \caption{Filtered mean trajectories by treatment group over visits.}
  \label{fig:filtered-mean-trajectories-by-treatment}
\end{figure}



\section{Estimation and Covariance Structure}

\subsection{Covariance structure}\label{subsec:covariance-structure}

As a starting point, The model formulation is specified as follows:
\begin{multline}
\label{eq:model1}
Y_{i} = I(TG_i = bud)\beta_{bud} + I(TG_i = ned)\beta_{ned} + I(TG_i = plbo)\beta_{plbo} \\
{}+ I(TG_i = bud)\,visit_{i}\beta_{visit \times bud}
+ I(TG_i = ned)\,visit_{i}\beta_{visit \times ned} \\
{}+ I(TG_i = plbo)\,visit_{i}\beta_{visit \times plbo}
+ age\_rz_{i}\,\beta_{age\_rz}
+ I(gender_{i} = m)\,\beta_{gender}\\
{}+ I(ethnic_{i} = b)\,\beta_{ethnic:black}
+ I(ethnic_{i} = h)\,\beta_{ethnic:hispanic}\\
{}+ I(ethnic_{i} = o)\,\beta_{ethnic:other})
+ \varepsilon_{i}\\
\varepsilon_{i} \sim N(0,\Sigma)
\end{multline}

The model included group-specific intercepts and
group-specific time slopes. The model adjusted for
three covariates: gender, ethnic, and age at randomization,
where female and white were set as reference levels.
The other covariates were not included, because they suffer
substantial missingness, which cannot be handled by
GLM.
The parameters are interpreted as follows:
\begin{itemize}
  \item $\beta_{bud}$: the baseline mean outcome for white, female subjects receiving budesonide, with baseline age set to 0.
  \item $\beta_{ned}$: the baseline mean outcome for white, female subjects receiving nedocromil, with baseline age set to 0.
  \item $\beta_{plbo}$: the baseline mean outcome for white, female subjects receiving placebo, with baseline age set to 0.
  \item $\beta_{visit \times bud}$: the change in mean outcome for each one-month increase in time among subjects receiving budesonide, after controlling for ethnicity, gender, and baseline age.
  \item $\beta_{visit \times ned}$: the change in mean outcome for each one-month increase in time among subjects receiving nedocromil, after controlling for ethnicity, gender, and baseline age.
  \item $\beta_{visit \times plbo}$: the change in mean outcome for each one-month increase in time among subjects receiving placebo, after controlling for ethnicity, gender, and baseline age.
  \item $\beta_{age\_rz}$: the change in mean outcome for one-year increase in baseline age between subjects, after controlling for treatment group, ethnicity, and gender.
  \item $\beta_{gender}$: the difference in mean outcome comparing male to female subjects, after controlling for treatment group, ethnicity, and baseline age.
  \item $\beta_{ethnic:black}$: the difference in mean outcome comparing black to white subjects, after controlling for treatment group, gender, and baseline age.
  \item $\beta_{ethnic:hispanic}$: the difference in mean outcome comparing hispanic to white subjects, after controlling for treatment group, gender, and baseline age.
  \item $\beta_{ethnic:other}$: the difference in mean outcome comparing subjects in other ethnicity group to white subjects, after controlling for treatment group, gender, and baseline age.
  \item $\varepsilon_i$: the random deviation of subject $i$'s observed outcomes from the mean outcomes.
\end{itemize}

\subsubsection{Correlation structure}

First, we fit the model assuming
an unstructured correlation matrix and equal variances over time in order to
determine the most appropriate correlation structure for the data.
However, there are 16 time points in the filtered data, which
means $16 \times 15/2 = 120$ parameters in the correlation matrix.
leading to prohibitive computational burden.
For computational efficiency, we selected a subset of time points
(0, 12, 24, 36, 48, 60, 72, 84, 96)
to fit the model using RMEL. R package \texttt{nlme} was used.
The estimated correlation matrix is shown in \cref{tab:fit-reml-rough-correlation}.

\begin{table}[htbp]
  \centering
  \caption{Estimated unstructured correlation matrix}
  \label{tab:fit-reml-rough-correlation}
  \small
\begin{tabular}{lrrrrrrrrr}
\hline
Visit & 0 & 12 & 24 & 36 & 48 & 60 & 72 & 84 & 96 \\
\hline
0 & 1.000 & 0.965 & 0.925 & 0.865 & 0.780 & 0.705 & 0.650 & 0.548 & 0.431 \\
12 & 0.965 & 1.000 & 0.957 & 0.905 & 0.816 & 0.726 & 0.660 & 0.556 & 0.424 \\
24 & 0.925 & 0.957 & 1.000 & 0.947 & 0.864 & 0.751 & 0.674 & 0.557 & 0.431 \\
36 & 0.865 & 0.905 & 0.947 & 1.000 & 0.933 & 0.823 & 0.729 & 0.596 & 0.465 \\
48 & 0.780 & 0.816 & 0.864 & 0.933 & 1.000 & 0.898 & 0.805 & 0.667 & 0.524 \\
60 & 0.705 & 0.726 & 0.751 & 0.823 & 0.898 & 1.000 & 0.936 & 0.806 & 0.665 \\
72 & 0.650 & 0.660 & 0.674 & 0.729 & 0.805 & 0.936 & 1.000 & 0.908 & 0.797 \\
84 & 0.548 & 0.556 & 0.557 & 0.596 & 0.667 & 0.806 & 0.908 & 1.000 & 0.917 \\
96 & 0.431 & 0.424 & 0.431 & 0.465 & 0.524 & 0.665 & 0.797 & 0.917 & 1.000 \\
\hline
\end{tabular}

\end{table}

Based on the pattern shown in \cref{tab:fit-reml-rough-correlation},
we assumed an exponential correlation structure,
since the times points are not equally spaced.
Although the correlation with fixed lag slightly decreases over time,
yet exponential correlation is still acceptable, which substantially reduces
the number of parameters to estimate.
Correlations depending on absolute time locations might imply
variance changing over time.
So next, we evaluated variance heterogeneity using a subset of
data with less time points (0, 24, 48, 72),
as we have more parameters to estimate. We fit the model
assuming an unstructured correlation matrix and unequal variances over time
using REML. The estimated correlation matrix is shown in
\cref{tab:fit-reml-rough2-correlation}, and the estimated
covariance matrix
is shown in \cref{tab:fit-reml-rough2-covariance}.

\begin{table}[htbp]
  \centering
  \caption{Estimated unstructured correlation matrix}
  \label{tab:fit-reml-rough2-correlation}
  \small
\begin{tabular}{lrrrr}
\hline
Visit & 0 & 24 & 48 & 72 \\
\hline
0 & 1.000 & 0.926 & 0.872 & 0.797 \\
24 & 0.926 & 1.000 & 0.904 & 0.778 \\
48 & 0.872 & 0.904 & 1.000 & 0.885 \\
72 & 0.797 & 0.778 & 0.885 & 1.000 \\
\hline
\end{tabular}

\end{table}

\begin{table}[htbp]
  \centering
  \caption{Estimated unstructured covariance matrix}
  \label{tab:fit-reml-rough2-covariance}
  \input{../hw3_code/output/fit_reml_rough2_covariance_matrix.tex}
\end{table}

As shown in \cref{tab:fit-reml-rough2-correlation}
and \cref{tab:fit-reml-rough2-covariance},
the model with unequal variances over time
modestly alleviates the issue of correlation depending
on absolute time locations. What is more important,
we find evidence of variance heterogeneity.

\subsubsection{Variance structure}

Then we evaluated different variance models.
We fit the model with exponetial correlation structure
to the data with 16 time points
using RMEL. The result is shown in \cref{tab:variance-model-compare}.
The model with time-specific variance exhibits the
smallest AIC. It is noteworthy that AIC is generally preferred
over BIC, as BIC's penalty grows with sample size
and tends to select parsimonious models.

\begin{table}[htbp]
  \centering
  \caption{Comparison of variance models by AIC and BIC}
  \label{tab:variance-model-compare}
  \begin{tabular}{lcrrr}
    \hline
    Variance model & Formulation & df & AIC & BIC \\
    \hline
    Time-specific variance & $\sigma^2(t) = \sigma^2_t$ & 28 & -1785.9 & -1585.0 \\
    Constant variance & $\sigma^2(t) = \sigma^2$ & 13 & -1699.5 & -1606.3 \\
    Power variance & $\sigma^2(t) = |t|^{2\theta}$ & 14 & -1708.3 & -1607.8 \\
    Constant plus power variance & $\sigma^2(t) = (\theta_1 + |t|^{\theta_2})^2$ & 15 & -1752.2 & -1644.6 \\
    Exponential variance & $\sigma^2(t) = exp(2\theta t) $ & 14 & -1751.3 & -1650.8 \\
    Constant plus proportion variance & $\sigma^2(t) = a^2 + b^2 t^2$ & 15 & -1727.6 & -1620.0 \\
    \hline
  \end{tabular}
\end{table}

\subsubsection{Baseline intercept}

Since the data is from a randomized trial,
it is reasonable to assume that the baseline means
are equal. Under this constraint, the model is specified
as follows:

\begin{multline}
\label{eq:model2}
Y_{i} = \beta_{0} + I(TG_i = bud)\,visit_{i}\beta_{visit \times bud}
+ I(TG_i = ned)\,visit_{i}\beta_{visit \times ned} \\
{}+ I(TG_i = plbo)\,visit_{i}\beta_{visit \times plbo}
+ age\_rz_{i}\,\beta_{age\_rz}
+ I(gender_{i} = m)\,\beta_{gender}\\
{}+ I(ethnic_{i} = b)\,\beta_{ethnic:black}
+ I(ethnic_{i} = h)\,\beta_{ethnic:hispanic}\\
{}+ I(ethnic_{i} = o)\,\beta_{ethnic:other})
+ \varepsilon_{i}\\
\varepsilon_{i} \sim N(0,\Sigma)
\end{multline}

We fit the two models with exponential correlation structure
and time-specific variance structure using ML,
and we performed a likelihood ratio test (LRT) to compare \cref{eq:model2}
with \cref{eq:model1}. The p-value is 0.271,
which suggests that the additional parameters in \cref{eq:model1}
do not significantly improve the model fit.
Moreover, the AIC of \cref{eq:model2}, -1864.3, is lower than
the AIC of \cref{eq:model1}, -1863.0. Thus we conclude that
\cref{eq:model2} appears to provide an adequate fit while
avoiding unnecessary complexity.

$beta_0$ is interpreted as baseline mean outcome for white, female
subjects with baseline age set to 0.

\subsubsection{Group-specific covariance matrices}

Lastly, We evaluated the model with group-specific covariance matrices.
Unfortunately, group-specific correlation matrix is not supported
by \texttt{nlme} package. Thus we evaluated the model with
group-specific variances against the model with homegeneous variances
across groups, while all other components were specified as in \cref{eq:model2},
and the models assume exponential correlation structure.

The p-value of LRT is $< 0.001$.
Additionally, the AIC of the more complex model is -1803.8,
lower than the AIC of the simpler model, -1795.3. The estimated
variances are shown in \cref{fig:variance-compare}. Thus we
conclude that the model
with group-specific variances provides a more adequate fit.
The final model formulation is as follows:


\begin{multline}
\label{eq:model3}
Y_{i} = \beta_{0} + I(TG_i = bud)\,visit_{i}\beta_{visit \times bud}
+ I(TG_i = ned)\,visit_{i}\beta_{visit \times ned} \\
{}+ I(TG_i = plbo)\,visit_{i}\beta_{visit \times plbo}
+ age\_rz_{i}\,\beta_{age\_rz}
+ I(gender_{i} = m)\,\beta_{gender}\\
{}+ I(ethnic_{i} = b)\,\beta_{ethnic:black}
+ I(ethnic_{i} = h)\,\beta_{ethnic:hispanic}\\
{}+ I(ethnic_{i} = o)\,\beta_{ethnic:other})
+ \varepsilon_{i}\\
\varepsilon_{i} \sim N(0,\Sigma_{TG_i})
\end{multline}

\begin{figure}[htbp]
   \centering
   \includegraphics[width=0.8\textwidth]{../hw3_code/output/variance_compare.png}
   \caption{Comparison of models with homogeneous and group-specific variances.}
   \label{fig:variance-compare}
\end{figure}
 
\subsubsection{Fitted mean trajectories}\label{subsubsec:fitted-mean-trajectories}
 
\begin{figure}[htbp]
    \centering
    \includegraphics[width=0.8\textwidth]{../hw3_code/output/fitted_observed_mean_compare.png}
    \caption{Observed and fitted mean trajectories by treatment group.}
    \label{fig:fitted-observed-mean-compare}
\end{figure}

To compute the
fitted mean for each treatment group at each time point,
we computed the mean of baseline age, the proportion of
each gender, and the proportion of each ethnic group.
The observed PostBD FEV1 means along with fitted PostBD FEV1 means
for each treatment group over time are illustrated in
\cref{fig:fitted-observed-mean-compare}. The model
shows an adequate fit to the data. However, the fitted
line underestimates the data. This is mainly due to a coincidence:
The difference in mean outcome among the three treatment groups
at later time points which have high leverage
is even smaller than that at baseline,
where the difference is solely due to randomness. Another
potential reason is that the correlation structure
we assume may be less than ideal, negatively affecting
the the mean model estimation.


\subsection{Estimation method}




\begin{table}[htbp]
  \caption{Coefficient estimates with model-based estimates}
  \label{tab:coef-model-based}
  \begingroup
  \catcode`\_=12\relax
  \pgfplotstabletypeset[
    col sep=comma,
    columns={term,Estimate,StdError,t,ddfm,p.val},
    columns/term/.style={column name=Term,string type},
    columns/Estimate/.style={column name=Estimate,fixed,precision=4},
    columns/StdError/.style={column name=Std. Error,fixed,precision=4},
    columns/t/.style={column name=t,fixed,precision=3},
    columns/ddfm/.style={column name=df,fixed,precision=0},
    columns/p.val/.style={column name=p-value,string type},
    every head row/.style={before row=\hline,after row=\hline},
    every last row/.style={after row=\hline}
  ]{../hw3_code/output/coef_model_based.csv}
  \endgroup
\end{table}

We fit the model using REML. The parameter estimates, model-based
standard errors, t-statistics, degress of freedom (Between-Within method), 
and p-value
are listed in \cref{tab:coef-model-based}.
The estimated covariance matrices for each treatment group
are reported in \cref{tab:covariance-matrix-budesonide},
\cref{tab:covariance-matrix-nedocromil} and
\cref{tab:covariance-matrix-placebo}.

\begin{table}[htbp]
  \centering
  \caption{Estimated covariance matrix under exponential correlation with unequal variance (Budesonide)}
  \label{tab:covariance-matrix-budesonide}
  \begingroup\let\small\scriptsize\makebox[\textwidth][c]{\small
\begin{tabular}{lrrrrrrrrrrrrrrrr}
\hline
Visit & 0 & 2 & 4 & 12 & 16 & 24 & 28 & 36 & 40 & 48 & 52 & 60 & 72 & 84 & 96 & 108 \\
\hline
0 & 0.383 & 0.396 & 0.381 & 0.355 & 0.334 & 0.321 & 0.329 & 0.322 & 0.314 & 0.316 & 0.296 & 0.287 & 0.254 & 0.230 & 0.210 & 0.196 \\
2 & 0.396 & 0.421 & 0.405 & 0.378 & 0.355 & 0.342 & 0.350 & 0.342 & 0.334 & 0.336 & 0.315 & 0.306 & 0.271 & 0.244 & 0.223 & 0.209 \\
4 & 0.381 & 0.405 & 0.402 & 0.375 & 0.352 & 0.339 & 0.347 & 0.340 & 0.332 & 0.334 & 0.313 & 0.303 & 0.269 & 0.243 & 0.222 & 0.207 \\
12 & 0.355 & 0.378 & 0.375 & 0.393 & 0.370 & 0.356 & 0.365 & 0.357 & 0.348 & 0.350 & 0.328 & 0.319 & 0.282 & 0.255 & 0.233 & 0.218 \\
16 & 0.334 & 0.355 & 0.352 & 0.370 & 0.369 & 0.356 & 0.364 & 0.356 & 0.348 & 0.350 & 0.328 & 0.318 & 0.282 & 0.254 & 0.232 & 0.217 \\
24 & 0.321 & 0.342 & 0.339 & 0.356 & 0.356 & 0.386 & 0.395 & 0.387 & 0.377 & 0.379 & 0.356 & 0.345 & 0.306 & 0.276 & 0.252 & 0.236 \\
28 & 0.329 & 0.350 & 0.347 & 0.365 & 0.364 & 0.395 & 0.429 & 0.420 & 0.410 & 0.412 & 0.387 & 0.375 & 0.332 & 0.300 & 0.274 & 0.256 \\
36 & 0.322 & 0.342 & 0.340 & 0.357 & 0.356 & 0.387 & 0.420 & 0.464 & 0.452 & 0.455 & 0.427 & 0.414 & 0.366 & 0.331 & 0.302 & 0.283 \\
40 & 0.314 & 0.334 & 0.332 & 0.348 & 0.348 & 0.377 & 0.410 & 0.452 & 0.468 & 0.471 & 0.442 & 0.428 & 0.379 & 0.342 & 0.313 & 0.293 \\
48 & 0.316 & 0.336 & 0.334 & 0.350 & 0.350 & 0.379 & 0.412 & 0.455 & 0.471 & 0.534 & 0.501 & 0.486 & 0.430 & 0.388 & 0.355 & 0.332 \\
52 & 0.296 & 0.315 & 0.313 & 0.328 & 0.328 & 0.356 & 0.387 & 0.427 & 0.442 & 0.501 & 0.499 & 0.484 & 0.428 & 0.387 & 0.353 & 0.330 \\
60 & 0.287 & 0.306 & 0.303 & 0.319 & 0.318 & 0.345 & 0.375 & 0.414 & 0.428 & 0.486 & 0.484 & 0.529 & 0.468 & 0.422 & 0.386 & 0.361 \\
72 & 0.254 & 0.271 & 0.269 & 0.282 & 0.282 & 0.306 & 0.332 & 0.366 & 0.379 & 0.430 & 0.428 & 0.468 & 0.496 & 0.448 & 0.409 & 0.383 \\
84 & 0.230 & 0.244 & 0.243 & 0.255 & 0.254 & 0.276 & 0.300 & 0.331 & 0.342 & 0.388 & 0.387 & 0.422 & 0.448 & 0.483 & 0.442 & 0.413 \\
96 & 0.210 & 0.223 & 0.222 & 0.233 & 0.232 & 0.252 & 0.274 & 0.302 & 0.313 & 0.355 & 0.353 & 0.386 & 0.409 & 0.442 & 0.483 & 0.452 \\
108 & 0.196 & 0.209 & 0.207 & 0.218 & 0.217 & 0.236 & 0.256 & 0.283 & 0.293 & 0.332 & 0.330 & 0.361 & 0.383 & 0.413 & 0.452 & 0.506 \\
\hline
\end{tabular}
}\endgroup
\end{table}

\begin{table}[htbp]
  \centering
  \caption{Estimated covariance matrix under exponential correlation with unequal variance (Nedocromil)}
  \label{tab:covariance-matrix-nedocromil}
  \begingroup\let\small\scriptsize\makebox[\textwidth][c]{\input{../hw3_code/output/fit_reml_c_exp_unequal_var_covariance_matrix_nedocromil.tex}}\endgroup
\end{table}

\begin{table}[htbp]
  \centering
  \caption{Estimated covariance matrix under exponential correlation with unequal variance (Placebo)}
  \label{tab:covariance-matrix-placebo}
  \begingroup\let\small\scriptsize\makebox[\textwidth][c]{\small
\begin{tabular}{lrrrrrrrrrrrrrrrr}
\hline
Visit & 0 & 2 & 4 & 12 & 16 & 24 & 28 & 36 & 40 & 48 & 52 & 60 & 72 & 84 & 96 & 108 \\
\hline
0 & 0.388 & 0.381 & 0.375 & 0.369 & 0.368 & 0.351 & 0.344 & 0.339 & 0.333 & 0.322 & 0.320 & 0.295 & 0.270 & 0.244 & 0.224 & 0.199 \\
2 & 0.381 & 0.387 & 0.380 & 0.374 & 0.373 & 0.356 & 0.349 & 0.344 & 0.337 & 0.327 & 0.325 & 0.299 & 0.274 & 0.247 & 0.227 & 0.202 \\
4 & 0.375 & 0.380 & 0.386 & 0.379 & 0.378 & 0.361 & 0.353 & 0.349 & 0.342 & 0.331 & 0.329 & 0.303 & 0.277 & 0.250 & 0.230 & 0.204 \\
12 & 0.369 & 0.374 & 0.379 & 0.419 & 0.418 & 0.399 & 0.391 & 0.386 & 0.378 & 0.367 & 0.365 & 0.336 & 0.307 & 0.277 & 0.255 & 0.226 \\
16 & 0.368 & 0.373 & 0.378 & 0.418 & 0.443 & 0.423 & 0.414 & 0.409 & 0.401 & 0.388 & 0.386 & 0.356 & 0.325 & 0.293 & 0.270 & 0.240 \\
24 & 0.351 & 0.356 & 0.361 & 0.399 & 0.423 & 0.455 & 0.446 & 0.440 & 0.431 & 0.418 & 0.416 & 0.383 & 0.350 & 0.316 & 0.290 & 0.258 \\
28 & 0.344 & 0.349 & 0.353 & 0.391 & 0.414 & 0.446 & 0.464 & 0.457 & 0.449 & 0.435 & 0.432 & 0.398 & 0.364 & 0.328 & 0.302 & 0.268 \\
36 & 0.339 & 0.344 & 0.349 & 0.386 & 0.409 & 0.440 & 0.457 & 0.508 & 0.499 & 0.483 & 0.480 & 0.442 & 0.404 & 0.365 & 0.336 & 0.298 \\
40 & 0.333 & 0.337 & 0.342 & 0.378 & 0.401 & 0.431 & 0.449 & 0.499 & 0.519 & 0.503 & 0.500 & 0.461 & 0.421 & 0.380 & 0.349 & 0.310 \\
48 & 0.322 & 0.327 & 0.331 & 0.367 & 0.388 & 0.418 & 0.435 & 0.483 & 0.503 & 0.549 & 0.546 & 0.503 & 0.460 & 0.415 & 0.382 & 0.339 \\
52 & 0.320 & 0.325 & 0.329 & 0.365 & 0.386 & 0.416 & 0.432 & 0.480 & 0.500 & 0.546 & 0.577 & 0.531 & 0.485 & 0.438 & 0.403 & 0.358 \\
60 & 0.295 & 0.299 & 0.303 & 0.336 & 0.356 & 0.383 & 0.398 & 0.442 & 0.461 & 0.503 & 0.531 & 0.551 & 0.504 & 0.455 & 0.418 & 0.371 \\
72 & 0.270 & 0.274 & 0.277 & 0.307 & 0.325 & 0.350 & 0.364 & 0.404 & 0.421 & 0.460 & 0.485 & 0.504 & 0.551 & 0.497 & 0.457 & 0.406 \\
84 & 0.244 & 0.247 & 0.250 & 0.277 & 0.293 & 0.316 & 0.328 & 0.365 & 0.380 & 0.415 & 0.438 & 0.455 & 0.497 & 0.537 & 0.494 & 0.438 \\
96 & 0.224 & 0.227 & 0.230 & 0.255 & 0.270 & 0.290 & 0.302 & 0.336 & 0.349 & 0.382 & 0.403 & 0.418 & 0.457 & 0.494 & 0.543 & 0.482 \\
108 & 0.199 & 0.202 & 0.204 & 0.226 & 0.240 & 0.258 & 0.268 & 0.298 & 0.310 & 0.339 & 0.358 & 0.371 & 0.406 & 0.438 & 0.482 & 0.512 \\
\hline
\end{tabular}
}\endgroup
\end{table}





\section{Variance Estimation }

\subsection{Comparison}

\begin{table}[htbp]
  \caption{Comparison of model-based and robust coefficient summaries}
  \label{tab:coef-model-vs-robust}
  \begin{tabular}{lrrrrl}
\hline
\multicolumn{6}{l}{Model-based} \\
\hline
Term & Estimate & Std. Error & t & df & p-value \\
\hline
(Intercept) & 1.7455 & 0.0894 & 19.530 & 689 & <0.001 \\
age\_rz & 0.0013 & 0.0096 & 0.139 & 689 & 0.445 \\
gendermale & 0.0263 & 0.0421 & 0.624 & 689 & 0.267 \\
ethnicblack & 0.0014 & 0.0626 & 0.023 & 689 & 0.491 \\
ethnichispanic & 0.0052 & 0.0703 & 0.074 & 689 & 0.471 \\
ethnicother & 0.0434 & 0.0745 & 0.583 & 689 & 0.280 \\
TGbudesonide:visitc & 0.0191 & 0.0004 & 43.168 & 8973 & <0.001 \\
TGnedocromil:visitc & 0.0199 & 0.0004 & 45.114 & 8973 & <0.001 \\
TGplacebo:visitc & 0.0187 & 0.0004 & 47.604 & 8973 & <0.001 \\
\hline
\multicolumn{6}{l}{Robust} \\
\hline
(Intercept) & 1.7455 & 0.0813 & 21.462 & 689 & <0.001 \\
age\_rz & 0.0013 & 0.0094 & 0.142 & 689 & 0.443 \\
gendermale & 0.0263 & 0.0405 & 0.649 & 689 & 0.258 \\
ethnicblack & 0.0014 & 0.0642 & 0.022 & 689 & 0.491 \\
ethnichispanic & 0.0052 & 0.0743 & 0.070 & 689 & 0.472 \\
ethnicother & 0.0434 & 0.0703 & 0.617 & 689 & 0.269 \\
TGbudesonide:visitc & 0.0191 & 0.0004 & 42.586 & 8973 & <0.001 \\
TGnedocromil:visitc & 0.0199 & 0.0004 & 45.183 & 8973 & <0.001 \\
TGplacebo:visitc & 0.0187 & 0.0004 & 48.995 & 8973 & <0.001 \\
\hline
\end{tabular}

\end{table}

The model-based and robust standard errors for
the fixed effects are presented in \cref{tab:coef-model-vs-robust}.
They are limitedly different, and this discrepancy does
not affect the statistical significance of the
fixed effects.


\subsection{Interpretation}

The model-based and robust standard errors
are broadly similar, indicating that the covariance model
was correctly specified.
Therefore, the model-based variance estimator
is more efficient.
Besides, the model was fit to a dataset with 695 subjects,
which provides large sample size,
so the robust variance estimator is also appropriate.
To summarize, these two variance estimators are comparable,
while the model-based one is slightly preferable.

\section{Hypothesis Testing}

\subsection{Primary hypothesis}

We performed a test of parallelism, with
$H_0: \beta_{visit \times bud} = \beta_{visit \times ned} = \beta_{visit \times plbo}$
versus $H_1: \text{not all three are equal}$.
The F statistic is 4.59, and the p-value is 0.010, which indicates
a significant result. Thus we reject the null hypothesis
and conclude that PostBD FEV1 changes differently over time
among the three treatment groups.

\subsection{Additional test}
We performed two tests to separately compare budesonide group versus
placebo group and nedocromil group versus placebo group.

First, we tested $H_0: \beta_{visit \times bud} = \beta_{visit \times plbo}$
versus $H_1: \beta_{visit \times bud} \neq \beta_{visit \times plbo}$.
The F statistic is 0.41, and the p-value is 0.521, which
indicates an insignficant result. Thus we do not reject the null
hypothesis and conclude that PostBD FEV1 does not
change differently over time between budesonide group and placebo
group.

Second, we tested $H_0: \beta_{visit \times ned} = \beta_{visit \times plbo}$
versus $H_1: \beta_{visit \times ned} \neq \beta_{visit \times plbo}$.
The F statistic is 4.48, and the p-value is 0.034, which suggests
a significant result. Thus we reject the null hypothesis and conclude
that PostBD FEV1 changes differently over time between nedocromil
group and placebo group.


\section{Summary and Interpretation}

We fit a general linear model using REML. We assumed an exponential
correlation structure, heterogeneouos variance over time and across
treatment groups, baseline mean outcomes constraint to be equal,
and linear outcome-time trend. These assumptions were chosen carefully
and have been justified using data visualization, LRT, or AIC.
We find that PostBD FEV1 changes differently over time among
the three treatment groups, where nedocromil group exhibits significant
difference relative to placebo group; however, budesonide shows
no signficant difference comparing to placebo group.
This implies that nedocromil is effective in improving
lung function among children with asthma.

\section{Code Availability}

The code is available at \url{https://github.com/CHENGHAO-WANG/767_homework/tree/main/hw3_code}.

\end{document}
